\chapter{CONCLUSION AND FUTURE ENHANCEMENTS}

\section{Conclusion}
The project titled Context-Aware Indoor Staff Presence Verification using BLE has been successfully designed and implemented. The system provides an efficient solution for tracking staff presence inside classrooms using Bluetooth Low Energy technology. The proposed system overcomes the limitations of traditional attendance systems such as manual registers, biometric systems, GPS-based tracking, and facial recognition methods. 

By using BLE technology, the system is able to work effectively in indoor environments where GPS fails. The use of ESP32 microcontrollers enables continuous detection of BLE signals, ensuring real-time monitoring of staff presence. One of the main advantages of the system is its context-aware approach. The system verifies both location and time before marking attendance, which improves accuracy and reduces false attendance entries. 

\section{Future Enhancements}
Although the system performs effectively, it can be further improved by adding additional features and enhancements. Future developments can focus on improving accuracy, scalability, and usability.
\begin{enumerate}
    \item \textbf{Multiple Receivers}: Accuracy can be enhanced by using trilateration with multiple BLE nodes. 
    \item \textbf{Mobile App}: Developing a native mobile application for push notification alerts. 
    \item \textbf{Behavioral Analytics}: Using AI to predict staff arrival patterns and optimize scheduling. 
\end{enumerate}
The system can also be extended to support multi-floor or large campus environments by integrating additional receivers and improving signal processing techniques.
